% Capitulo: A chegada da IA
\chapter{A chegada da IA}
\label{cap:ia}

\rascunho{
  \textbf{Ideia do capítulo.} O que a IA generativa mudou no trabalho com dados e o que ela não mudou.
  \begin{itemize}
    \item Modelos de linguagem como copilotos: escrever SQL, explicar código, documentar.
    \item Text-to-SQL e a promessa de \enquote{conversar com os dados}: onde funciona e onde engana.
    \item Agentes que executam tarefas inteiras e a nova pergunta: quem revisa?
    \item O que não mudou: definição de métrica, qualidade do dado, contexto de negócio.
    \item O que a IA passou a pedir dos dados (treino e contexto) fica para a
      Parte~\ref{parte:modelagem}; o custo em tokens, para o capítulo~\ref{cap:custo-consumo}.
  \end{itemize}
}

\section{O copiloto na mesa do analista}
\label{sec:copiloto}

\pendente{escrever}

\section{Conversar com os dados}
\label{sec:text-to-sql}

\pendente{escrever}

\section{Agentes e a pergunta de quem revisa}
\label{sec:agentes}

\pendente{escrever}

\section{O que continua igual}
\label{sec:continua-igual}

\pendente{escrever}
